\scanpage{175}
% archive sheet number: 174
\deleted{Supposons $A_0$ \unsure{fblt} dense dans $\Omega'$, i.e. $A_0$ en dualité avec $\Omega$}
Soit $K$ la classe des $\Omega$-modules $V$ \inserted{de dim. finie sur $k$} tels que
pour $x\in V$, $x'\in V'$, $f_{xx'}$ soit dans $A_0$. Cette
classe est isomorphe à la classe des $A_0$-comodules
(un tel module définit un
comodule sans hyp. sur la dimension ;
la réciproque marche, grâce à la dimension
finie). C'est donc aussi la classe des
$\widehat\Omega$-modules, où $\widehat\Omega$ est le complété
\deleted{\unsure{faible}} de $\Omega$ pour $\sigma(\Omega,A_0)$.

\unsure{On va interpréter|Or on interprète} autrement la donnée
d'un $A_0$ comme ci-dessus. $A_0$ définit, \inserted{on l'a vu,}
un ``groupe de représentations'' de dim finie
de $\Omega$ (i.e. classe de représentations stable
\unsure{par|sous} équivalence, passage aux sous-trucs et
quotients, sommes directes). Supposons
donné réciproquement un tel groupe $\unsure{\mathscr G|\mathcal G|G}$.
Soit $A_0$ l'ensemble des coefficients de
\unsure{ttes} les représentations de $\mathscr G$. \deleted{Alors} Je dis
que $A_0$ satisfait aux conditions précédentes :
en effet, \inserted{\textsuperscript{\unsure{(2)|(1)}}} si $f$ est un coéfficient d'une
représentation de dim finie de $\Omega$ sur $V$,
on a vu que \deleted{$\delta f=\sum f_i\otimes f'_i$, les}
\deleted{$f_i$ et $f'_i$ coefficients de la \unsure{même} représentation.}
\inserted{\unsure{les sommes de coefficients|les sommes de coéfficients} de la}
représentation forment déjà une coalgèbre
